% GENERATED FILE. Edit canonical lessons, then run npm run generate:books.
% canonical-source-hash: fnv1a64:32ec0017b03cdfc3
% canonical-lessons: BN-C473-songkromon, BN-C473-jibanu, BN-C473-chhoyache, BN-C473-mohamari, BN-C473-akranto-hooya

\chapter{Germs and Infection}
\label{ch:bn-germs-and-infection}
\hlchaptermodality{473}

\begin{chapteropening}
\textbf{By the end of this chapter:} \emph{I can talk about germs, infection and a disease spreading.}

Complete the last lesson of chapter 473: I can talk about germs, infection and a disease spreading.
\end{chapteropening}

\section[\texorpdfstring{\bn{সংক্রমণ} (sôngkrômôn)}{সংক্রমণ (sôngkrômôn)}]{\bn{সংক্রমণ} (sôngkrômôn) — an infection}

\label{lesson:BN-C473-songkromon}



\begin{quote}
Before the new one: say the Bengali for a vaccine, then the Bengali for to recover.
\end{quote}

\subsection*{You'll want to know: \bn{সংক্রমণ}}
\textbf{\bn{সংক্রমণ}} — \emph{sôngkrômôn} — \textquotedblleft{}an infection\textquotedblright{}.

\begin{quote}
\bn{কাটা} \bn{জায়গায়} \bn{সংক্রমণ} \bn{হলে} \bn{জ্বর} \bn{আসে।}
\end{quote}

\begin{quote}
\emph{kāṭā jāygāy sôngkrômôn hôle jbôr āse.}
\end{quote}

\begin{quote}
If a cut gets infected, a fever follows.
\end{quote}

\begin{cousinweb}
\textbf{\bn{সংক্রমণ}} (\emph{sôngkrômôn}) is a tatsama, Sanskrit \emph{saṃkramaṇa}, \textquotedblleft{}going together, transference\textquotedblright{} (Monier-Williams); Bengali uses it for a disease passing from one body to another.
\end{cousinweb}

\begin{grammarlens}[title={What you've built}]
One more word for describing what you lived and felt.
\end{grammarlens}

\subsection*{Your turn}
\begin{itemize}
\raggedright
  \item \emph{Say it:} \emph{sôngkrômôn}
  \item \emph{Say it:} \emph{sôngkrômôn}, once more
  \item \emph{Say it:} \emph{kāṭā jāygāy sôngkrômôn hôle jbôr āse.} — If a cut gets infected, a fever follows.
  \item \emph{Recall:} say the Bengali for a vaccine, then the Bengali for to recover, then say \emph{sôngkrômôn} again
\end{itemize}

\begin{tcolorbox}[breakable,colback=teal!4,colframe=teal!35!black,title={Before you move on}]
What does \bn{সংক্রমণ} mean? (\textquotedblleft{}An infection\textquotedblright{}.) Say it once more.
\end{tcolorbox}
\section[\texorpdfstring{\bn{জীবাণু} (jībānu)}{জীবাণু (jībānu)}]{\bn{জীবাণু} (jībānu) — a germ}

\label{lesson:BN-C473-jibanu}



\begin{quote}
Before the new one: say the Bengali for to recover, then the Bengali for an infection.
\end{quote}

\subsection*{You'll want to know: \bn{জীবাণু}}
\textbf{\bn{জীবাণু}} — \emph{jībānu} — \textquotedblleft{}a germ\textquotedblright{}.

\begin{quote}
\bn{সাবান} \bn{দিয়ে} \bn{হাত} \bn{ধুলে} \bn{জীবাণু} \bn{মরে।}
\end{quote}

\begin{quote}
\emph{sābān diye hāt dhule jībānu môre.}
\end{quote}

\begin{quote}
Washing your hands with soap kills germs.
\end{quote}

\begin{cousinweb}
\textbf{\bn{জীবাণু}} (\emph{jībānu}) is a modern formation from Sanskrit \emph{jīva}, \textquotedblleft{}living\textquotedblright{}, and \emph{aṇu}, \textquotedblleft{}minute, an atom\textquotedblright{} (Monier-Williams): a tiny living thing.
\end{cousinweb}

\begin{grammarlens}[title={What you've built}]
One more word for describing what you lived and felt.
\end{grammarlens}

\subsection*{Your turn}
\begin{itemize}
\raggedright
  \item \emph{Say it:} \emph{jībānu}
  \item \emph{Say it:} \emph{jībānu}, once more
  \item \emph{Say it:} \emph{sābān diye hāt dhule jībānu môre.} — Washing your hands with soap kills germs.
  \item \emph{Recall:} say the Bengali for to recover, then the Bengali for an infection, then say \emph{jībānu} again
\end{itemize}

\begin{tcolorbox}[breakable,colback=teal!4,colframe=teal!35!black,title={Before you move on}]
What does \bn{জীবাণু} mean? (\textquotedblleft{}A germ\textquotedblright{}.) Say it once more.
\end{tcolorbox}
\section[\texorpdfstring{\bn{ছোঁয়াচে} (chhõyāche)}{ছোঁয়াচে (chhõyāche)}]{\bn{ছোঁয়াচে} (chhõyāche) — contagious; catching}

\label{lesson:BN-C473-chhoyache}



\begin{quote}
Before the new one: say the Bengali for an infection, then the Bengali for a germ.
\end{quote}

\subsection*{You'll want to know: \bn{ছোঁয়াচে}}
\textbf{\bn{ছোঁয়াচে}} — \emph{chhõyāche} — \textquotedblleft{}contagious; catching\textquotedblright{}.

\begin{quote}
\bn{এই} \bn{রোগটা} \bn{ছোঁয়াচে}, \bn{তাই} \bn{রোগীকে} \bn{আলাদা} \bn{রাখতে} \bn{হবে।}
\end{quote}

\begin{quote}
\emph{ei rogṭā chhõyāche, tāi rogīke ālādā rākhte hôbe.}
\end{quote}

\begin{quote}
This disease is contagious, so the patient must be kept apart.
\end{quote}

\begin{cousinweb}
\textbf{\bn{ছোঁয়াচে}} (\emph{chhõyāche}) is built on \emph{chhõyā}, \textquotedblleft{}to touch\textquotedblright{}, which CDIAL lists for Bengali under Sanskrit \emph{chupati}, \textquotedblleft{}touches\textquotedblright{} (CDIAL 5055), with an adjective ending: caught by touching.
\end{cousinweb}

\begin{grammarlens}[title={What you've built}]
One more word for describing what you lived and felt.
\end{grammarlens}

\subsection*{Your turn}
\begin{itemize}
\raggedright
  \item \emph{Say it:} \emph{chhõyāche}
  \item \emph{Say it:} \emph{chhõyāche}, once more
  \item \emph{Say it:} \emph{ei rogṭā chhõyāche, tāi rogīke ālādā rākhte hôbe.} — This disease is contagious, so the patient must be kept apart.
  \item \emph{Recall:} say the Bengali for an infection, then the Bengali for a germ, then say \emph{chhõyāche} again
\end{itemize}

\begin{tcolorbox}[breakable,colback=teal!4,colframe=teal!35!black,title={Before you move on}]
What does \bn{ছোঁয়াচে} mean? (\textquotedblleft{}Contagious; catching\textquotedblright{}.) Say it once more.
\end{tcolorbox}
\section[\texorpdfstring{\bn{মহামারি} (môhāmāri)}{মহামারি (môhāmāri)}]{\bn{মহামারি} (môhāmāri) — an epidemic}

\label{lesson:BN-C473-mohamari}



\begin{quote}
Before the new one: say the Bengali for a germ, then the Bengali for contagious.
\end{quote}

\subsection*{You'll want to know: \bn{মহামারি}}
\textbf{\bn{মহামারি}} — \emph{môhāmāri} — \textquotedblleft{}an epidemic\textquotedblright{}.

\begin{quote}
\bn{একশো} \bn{বছর} \bn{আগে} \bn{শহরে} \bn{এক} \bn{মহামারি} \bn{হল।}
\end{quote}

\begin{quote}
\emph{eksho bôchhôr āge shôhôre ek môhāmāri hôlô.}
\end{quote}

\begin{quote}
A hundred years ago an epidemic broke out in the city.
\end{quote}

\begin{cousinweb}
\textbf{\bn{মহামারি}} (\emph{môhāmāri}) is a tatsama, Sanskrit \emph{mahāmārī}, \textquotedblleft{}the great destroyer\textquotedblright{}, a pestilence (Monier-Williams).
\end{cousinweb}

\begin{grammarlens}[title={What you've built}]
One more word for describing what you lived and felt.
\end{grammarlens}

\subsection*{Your turn}
\begin{itemize}
\raggedright
  \item \emph{Say it:} \emph{môhāmāri}
  \item \emph{Say it:} \emph{môhāmāri}, once more
  \item \emph{Say it:} \emph{eksho bôchhôr āge shôhôre ek môhāmāri hôlô.} — A hundred years ago an epidemic broke out in the city.
  \item \emph{Recall:} say the Bengali for a germ, then the Bengali for contagious, then say \emph{môhāmāri} again
\end{itemize}

\begin{tcolorbox}[breakable,colback=teal!4,colframe=teal!35!black,title={Before you move on}]
What does \bn{মহামারি} mean? (\textquotedblleft{}An epidemic\textquotedblright{}.) Say it once more.
\end{tcolorbox}
\section[\texorpdfstring{\bn{আক্রান্ত} \bn{হওয়া} (ākrāntô hôoyā)}{আক্রান্ত হওয়া (ākrāntô hôoyā)}]{\bn{আক্রান্ত} \bn{হওয়া} (ākrāntô hôoyā) — to be struck down (by an illness); to be attacked}

\label{lesson:BN-C473-akranto-hooya}



\begin{quote}
Before the new one: say the Bengali for contagious, then the Bengali for an epidemic.
\end{quote}

\subsection*{You'll want to know: \bn{আক্রান্ত} \bn{হওয়া}}
\textbf{\bn{আক্রান্ত} \bn{হওয়া}} — \emph{ākrāntô hôoyā} — \textquotedblleft{}to be struck down (by an illness); to be attacked\textquotedblright{}.

\begin{quote}
\bn{বন্যার} \bn{পরে} \bn{গ্রামের} \bn{অনেক} \bn{লোক} \bn{জ্বরে} \bn{আক্রান্ত} \bn{হল।}
\end{quote}

\begin{quote}
\emph{bônyār pôre grāmer ônek lok jbôre ākrāntô hôlô.}
\end{quote}

\begin{quote}
After the flood many people in the village went down with fever.
\end{quote}

Past (he/she …): \bn{তিনি} \bn{জ্বরে} \bn{আক্রান্ত} \bn{হলেন}, \emph{tini jbôre ākrāntô hôlen}.

The illness takes -\bn{এ}: \bn{জ্বরে} \bn{আক্রান্ত}, down with fever.

\begin{cousinweb}
\textbf{\bn{আক্রান্ত} \bn{হওয়া}} (\emph{ākrāntô hôoyā}) is Sanskrit \emph{ākrānta}, \textquotedblleft{}overcome, overrun, attacked\textquotedblright{} (Monier-Williams), with \emph{hôoyā}, \textquotedblleft{}to be, to become\textquotedblright{} (CDIAL 9416): to be overrun, here by a disease.
\end{cousinweb}

\begin{grammarlens}[title={What you've built}]
One more word for describing what you lived and felt.
\end{grammarlens}

\subsection*{Your turn}
\begin{itemize}
\raggedright
  \item \emph{Say it:} \emph{ākrāntô hôoyā}
  \item \emph{Say it:} \emph{ākrāntô hôoyā}, once more
  \item \emph{Say it:} \emph{bônyār pôre grāmer ônek lok jbôre ākrāntô hôlô.} — After the flood many people in the village went down with fever.
  \item \emph{Recall:} say the Bengali for contagious, then the Bengali for an epidemic, then say \emph{ākrāntô hôoyā} again
\end{itemize}

\begin{tcolorbox}[breakable,colback=teal!4,colframe=teal!35!black,title={Before you move on}]
What does \bn{আক্রান্ত} \bn{হওয়া} mean? (\textquotedblleft{}To be struck down (by an illness); to be attacked\textquotedblright{}.) Say it once more.
\end{tcolorbox}
